\documentclass[12pt]{article}
\usepackage[margin=1in]{geometry}
\usepackage{booktabs,threeparttable}

\begin{document}

\providecommand{\StructuralRobustnessPath}{appendix/structural-robustness}

\subsection{Asymmetric Classification and Regularization of the Reporting Process}
\label{app:asymmetric-reporting-priors}

The baseline structural model maps definite classifications of recognized
peers directly into observability. As a robustness exercise, we instead
estimate the complete reporting process using peer reports linked to
administrative take-up. Conditional on recognizing a peer, the probabilities
of reporting Yes, No, or Don't Know depend on the peer's true take-up status,
treatment assignment, distance, and treatment-by-distance interactions. The
resulting truth-specific response matrix, including false-positive and
false-negative classifications, enters the equilibrium fixed point.

The treatment-specific report-by-distance interactions are estimated from a
relatively small number of treatment-distance cells. We report three ways of
handling this weakly informed portion of the reporting channel. The
unrestricted specification assigns the twelve treatment-specific distance
contrasts independent $N(0,0.25)$ priors. The tighter-prior specification
changes these to $N(0,0.10)$. A design-pooled specification instead writes
each slope as a No-Signal or Any-Signal group slope plus a centered
within-group deviation: Control and Calendar deviations sum to zero, as do Ink
and Bracelet deviations. A common half-normal scale shrinks those within-group
deviations while allowing the two experimentally meaningful signal groups to
have different reporting gradients. All other priors and likelihood
components are unchanged.

Table~\ref{tab:tight-asymmetric-report-multipliers} reports both local point
multipliers and a finite-change multiplier over the 500--2500m interval. The
tighter prior yields multipliers close to one for Control, Calendar, and Ink.
For Bracelet, the local multiplier declines with distance and has a posterior
median of $-0.12$ at 2500m, but its 95 percent credible interval includes zero
widely. The corresponding finite-change multiplier is $0.34$, with a 95
percent credible interval of $(-0.04,0.82)$. Thus the specification indicates
substantial attenuation of the direct access-cost effect for Bracelet, while
providing limited evidence that the response reverses over the full distance
interval.

The structural average treatment effects remain stable under this
regularization. The combined effects are $-2.5$ percentage points for Ink,
$+4.3$ percentage points for Calendar, and $+7.2$ percentage points for
Bracelet. The sensitivity is therefore concentrated in the decomposition of
the distance response rather than in recovery of the average treatment
effects.

Because the four multiplier estimates are computed from the same posterior
draw, separate arm-by-arm intervals discard useful covariance. Our primary
robustness summary therefore compares the average multiplier for Control and
Calendar (No Signal) with the average for Ink and Bracelet (Any Signal) within
each draw. Table~\ref{tab:asymmetric-report-contrasts} shows that this paired
contrast is similar across the unrestricted, hierarchical, tighter-prior, and
design-pooled reporting models. Under design pooling, the contrast is 0.47
with a 95 percent interval of $(0.18,0.74)$ at 1500m and 0.45
$(0.17,0.71)$ for the finite 500--2500m change. The corresponding within-group
shrinkage scale has posterior median 0.12 and 95 percent interval
$(0.03,0.27)$ under its half-normal scale-0.25 prior. Sampling diagnostics are
clean (maximum $\widehat R=1.003$, minimum bulk ESS 843, no divergences, and no
maximum-treedepth transitions). A Newton-first one-dimensional GQ solver
reproduces the original algebra-solver multiplier values exactly on every draw
where the latter converges and returns finite multiplier values for all 3,200
posterior draws.

The evidence is not uniform over the full distance function. At 500m the
design-pooled contrast is 0.03, with a 95 percent interval from $-0.23$ to
$0.27$. The minimum over the three reported distances is also 0.03, with an
interval from $-0.23$ to $0.26$. We therefore interpret the exercise
as supporting a positive pooled signal-versus-no-signal contrast at the
paper's 1500m reference distance and over the observed 500--2500m finite
change, not as establishing that every arm-specific multiplier differs from
one at every distance.

\begin{table}[htbp]
  \centering
  \small
  \caption{Social multipliers with asymmetric classification and tighter priors}
  \label{tab:tight-asymmetric-report-multipliers}
  \begin{threeparttable}
    \begin{tabular}{lcccc}
\toprule
& \multicolumn{3}{c}{Point multiplier} & Finite change \\
\cmidrule(lr){2-4}\cmidrule(l){5-5}
Treatment & 500m & 1500m & 2500m & 500--2500m \\
\midrule
Bracelet & 0.85 & 0.34 & -0.12 & 0.34 \\
 & (0.67, 1.09) & (-0.04, 0.81) & (-0.94, 0.78) & (-0.04, 0.82) \\
Calendar & 0.93 & 0.93 & 0.97 & 0.93 \\
 & (0.83, 1.11) & (0.72, 1.09) & (0.69, 1.08) & (0.74, 1.09) \\
Ink & 0.89 & 0.86 & 0.92 & 0.87 \\
 & (0.75, 1.00) & (0.66, 1.01) & (0.67, 1.03) & (0.68, 1.01) \\
Control & 0.95 & 1.05 & 1.14 & 1.05 \\
 & (0.78, 1.20) & (0.84, 1.29) & (0.97, 1.33) & (0.86, 1.26) \\
\bottomrule
\end{tabular}

    \begin{tablenotes}\footnotesize
      \item \textit{Notes:} Entries are posterior medians; parentheses contain
      95 percent posterior credible intervals. The first three columns report
      local point multipliers at the indicated distances. The final column
      reports the finite change in the equilibrium cutoff between 500m and
      2500m, normalized by the corresponding direct access-cost change. A
      multiplier of one indicates no amplification, a value between zero and
      one indicates attenuation, and a negative value indicates a reversal of
      the local comparative static. The twelve treatment-specific
      report-by-distance contrasts have independent $N(0,0.10)$ priors. All
      other priors and likelihood components are unchanged.
    \end{tablenotes}
  \end{threeparttable}
\end{table}

\begin{table}[htbp]
  \centering
  \scriptsize
  \caption{Paired multiplier contrast across asymmetric-reporting specifications}
  \label{tab:asymmetric-report-contrasts}
  {\setlength{\tabcolsep}{4pt}%
  \begin{tabular}{lccccc}
\toprule
Specification & 500m & 1500m & 2500m & Finite & Grid minimum \\
\midrule
Unrestricted multinomial, N(0, 0.25) & -0.02 & 0.51 & 0.87 & 0.49 & -0.02 \\
 & (-0.27, 0.22) & (0.26, 0.78) & (0.36, 1.31) & (0.24, 0.73) & (-0.27, 0.22) \\
Hierarchical multinomial, half-N(0, 0.25) scales & -0.04 & 0.52 & 0.86 & 0.48 & -0.04 \\
 & (-0.31, 0.21) & (0.25, 0.79) & (0.35, 1.31) & (0.23, 0.73) & (-0.31, 0.21) \\
Tighter multinomial, N(0, 0.10) & 0.08 & 0.39 & 0.65 & 0.38 & 0.08 \\
 & (-0.09, 0.27) & (0.11, 0.65) & (0.16, 1.14) & (0.10, 0.64) & (-0.09, 0.27) \\
Design-pooled multinomial & 0.03 & 0.47 & 0.79 & 0.45 & 0.03 \\
 & (-0.23, 0.27) & (0.18, 0.74) & (0.23, 1.30) & (0.17, 0.71) & (-0.23, 0.26) \\
\bottomrule
\end{tabular}
}
  \begin{minipage}{0.94\linewidth}
    \vspace{0.4em}\footnotesize
    \textit{Notes:} Entries are posterior medians; parentheses contain 95
    percent credible intervals. Each entry is the within-draw difference
    between the average multiplier for Control and Calendar (No Signal) and
    the average for Ink and Bracelet (Any Signal). ``Finite'' uses the change
    from 500m to 2500m. ``Grid minimum'' takes the minimum paired contrast over
    the three displayed distances within each posterior draw before computing
    summaries. No intervals are truncated and no confidence level is changed.
  \end{minipage}
\end{table}

\end{document}
